% Libro 4 · Mundo 3 · Fuentes y voces (app: pu-fuentes)
\mundo{Fuentes y voces}{pu-fuentes}{Un texto rara vez habla solo: cita expertos, recoge testimonios, presenta posturas
contrarias para responderlas. Saber quién dice qué, con qué autoridad y para qué lo trae el autor es clave para no
confundir la opinión del autor con la de sus citados, y para evaluar si una fuente es confiable.}{%
\item \textbf{3.1} Citas y voces dentro del texto
\item \textbf{3.2} Posturas y citas externas}

\leccion{pu-fu-1}{Citas y voces}{Juego de asignar cada idea a la voz que la sostiene.}

\tecnica{La técnica: ¿para qué trae el autor esta voz?}
Cuando un autor cita a alguien, puede hacerlo por tres razones: para \clave{apoyarse} (cita a un experto que dice lo
mismo), para \clave{oponerse} (presenta una idea que luego refuta) o para \clave{ilustrar} (un testimonio que vuelve
concreta la idea). Las marcas ayudan: «como señala…» suele apoyar; «hay quienes creen…», «se suele decir…» suele preparar
una refutación.

\figura{paraque}{Las tres razones por las que un autor trae una cita.}

\begin{regla}
\begin{itemize}
  \item Distingue la voz del autor de las voces citadas.
  \item Pregunta para qué está la cita: apoyar, oponer o ilustrar.
  \item «Hay quienes», «se dice», «algunos piensan» casi siempre presentan una idea que el autor no comparte.
\end{itemize}
\end{regla}

\begin{texto}{¿Leemos peor en pantalla?}{Artículo de opinión · texto original escrito para este libro}
\pa Hay quienes sostienen que el soporte da igual: un texto es un texto, en papel o en pantalla. La comodidad de los
dispositivos, dicen, compensa cualquier diferencia. Es una idea razonable, pero la evidencia acumulada en los últimos
años invita a matizarla.

\pa Varias investigaciones con estudiantes han encontrado que, en textos largos y complejos, la comprensión tiende a ser
algo mejor en papel. Una explicación posible es que en pantalla leemos con la costumbre de las redes: rápido, saltando,
buscando palabras clave. «En pantalla no leemos: escaneamos», resume una profesora de lenguaje que entrevisté.

\pa Esto no convierte a las pantallas en enemigas. Para consultar datos, buscar información o leer textos breves, son
herramientas excelentes. El problema aparece cuando usamos el mismo modo de lectura para todo.

\pa Quizás la conclusión sea que no hay que elegir entre papel y pantalla, sino aprender a cambiar de velocidad: saber
cuándo escanear y cuándo detenerse.
\end{texto}

\subsection{Preguntas comentadas}
\begin{comentada}{La frase «En pantalla no leemos: escaneamos» \pf{2}}
\mal{A}{es la conclusión del autor.}{La conclusión está en el \S\,4.}
\bien{B}{es la voz de una profesora citada para ilustrar la explicación del autor.}{Apoya e ilustra lo dicho antes sobre
la lectura «rápido, saltando».}
\mal{C}{expresa la idea que el autor refuta.}{\emph{Dice lo contrario}: el autor la usa a su favor.}
\mal{D}{es un dato estadístico.}{Es un testimonio, no una cifra.}
\end{comentada}

\begin{reto}
Un autor escribe: «Hay quienes creen que leer en pantalla es igual que leer en papel. Los estudios, sin embargo, sugieren
algo distinto». ¿Cuál es la postura del autor?
\opcionesl{Leer en pantalla es igual que en papel.}{Probablemente no es igual leer en pantalla que en papel.}{Los estudios son falsos.}{No tiene postura.}
\end{reto}

\begin{retoresuelto}
«Hay quienes creen…» presenta la idea ajena; «sin embargo» introduce la del autor: no es igual. \textbf{Opción B.}
\end{retoresuelto}

\taller
\piso{Como en el examen}{opción múltiple}
\begin{ejercicios}
\item La idea de que «el soporte da igual» \pf{1} es
\opcionesl{la tesis del autor.}{una postura que el autor presenta para luego matizarla.}{la opinión de la profesora.}{un resultado de las investigaciones.}
\sol{B.}
\item El autor usa las investigaciones del \S\,2 para
\opcionesd{apoyar su postura}{refutarse a sí mismo}{entretener}{cambiar de tema}
\sol{A.}
\item Según el autor, las pantallas
\opcionesl{deben prohibirse en los colegios.}{son excelentes para ciertas tareas, como consultar datos.}{siempre dan mejor comprensión.}{solo sirven para redes sociales.}
\sol{B. \S\,3.}
\end{ejercicios}

\bitacora

\leccion{pu-fu-2}{Posturas y citas externas}{Juego de clasificar fuentes según su confiabilidad.}

\tecnica{La técnica: tres preguntas a una fuente}
Cuando un texto se apoya en una fuente, conviene preguntar:
\begin{enumerate}
  \item \textbf{¿Quién lo dice?} ¿Tiene conocimiento del tema?
  \item \textbf{¿Qué interés tiene?} ¿Gana algo si le creemos?
  \item \textbf{¿Cómo lo sabe?} ¿Es un estudio, una experiencia personal, un rumor?
\end{enumerate}
Una fuente no es falsa por tener interés, ni verdadera por tener un título; pero estas preguntas permiten pesar cuánto
apoyo da a una afirmación.

\subsubsection{Posturas}
Ante un tema, un texto puede mostrar varias \clave{posturas}: a favor, en contra, intermedia. Las preguntas piden
identificar cuál es la del autor y cuáles son las de otros, y a veces imaginar qué diría cada uno.

\figura{fuente}{Tres preguntas para pesar una fuente.}

\begin{regla}
\begin{itemize}
  \item Quién, qué interés, cómo lo sabe.
  \item Experiencia personal: sirve para ilustrar, no para generalizar.
  \item Separa la postura del autor de las que solo reporta.
\end{itemize}
\end{regla}

\begin{texto}{¿Uniforme sí o no?}{Artículo · texto original escrito para este libro (personas imaginarias)}
\pa El consejo directivo de un colegio discute si eliminar el uniforme. La rectora, Marta Gil, está a favor de mantenerlo:
«Evita la competencia por la ropa de marca y hace que todos se sientan parte de algo».

\pa Un grupo de estudiantes piensa lo contrario. Para Sebastián, del grado once, «el uniforme no borra las diferencias:
todos sabemos quién tiene los tenis más caros». Propone que cada estudiante se vista como quiera, con algunas reglas
básicas.

\pa La asociación de padres, por su parte, advierte un tema práctico: el uniforme es costoso al comienzo del año, pero
resulta más barato que comprar ropa variada. Una empresa de uniformes que surte al colegio ha enviado una carta defendiendo
su uso «por razones pedagógicas».

\pa El consejo decidirá el próximo mes. Mientras tanto, una encuesta interna muestra que los estudiantes están divididos
casi por la mitad.
\end{texto}

\subsection{Preguntas comentadas}
\begin{comentada}{¿Cuál de las fuentes citadas tiene un interés económico más directo en el resultado?}
\mal{A}{La rectora.}{Tiene una postura, pero el texto no le atribuye interés económico.}
\mal{B}{Sebastián.}{\emph{No lo dice}.}
\mal{C}{La asociación de padres.}{Su argumento es económico, pero se refiere a ahorrar, no a ganar.}
\bien{D}{La empresa de uniformes.}{Gana si el uniforme se mantiene; por eso sus «razones pedagógicas» merecen cautela.}
\end{comentada}

\begin{reto}
Para respaldar que una bebida «mejora la memoria», ¿cuál fuente sería más confiable?
\opcionesl{El anuncio de la empresa que la vende.}{Un estudio publicado por investigadores independientes, con grupo de comparación.}{El comentario de un usuario en una red social.}{Un amigo que la toma todos los días.}
\end{reto}

\begin{retoresuelto}
El estudio independiente con grupo de comparación responde mejor a las tres preguntas: conocimiento, sin interés directo,
método. \textbf{Opción B.}
\end{retoresuelto}

\taller
\piso{Como en el examen}{opción múltiple}
\begin{ejercicios}
\item ¿Cuál es la postura del autor del artículo?
\opcionesl{Está a favor del uniforme.}{Está en contra del uniforme.}{No expresa una postura; presenta varias.}{Apoya a la empresa.}
\sol{C.}
\item El argumento de Sebastián responde a la idea de la rectora de que el uniforme
\opcionesl{es barato.}{evita la competencia por la ropa de marca.}{es pedagógico.}{es incómodo.}
\sol{B. «No borra las diferencias».}
\item ¿Qué diría probablemente la asociación de padres si se eliminara el uniforme?
\opcionesl{Que los gastos en ropa podrían aumentar.}{Que los estudiantes estudiarán más.}{Que la empresa debe cerrar.}{Que la rectora tenía razón en todo.}
\sol{A.}
\end{ejercicios}

\bitacora

\evaluacion
\begin{texto}{La siesta en el trabajo}{Artículo · texto original escrito para este libro (personas imaginarias)}
\pa Algunas empresas han empezado a permitir siestas cortas a sus empleados. Para el médico del sueño Andrés Vélez, «una
siesta de veinte minutos mejora la atención de la tarde más que un café». Advierte, eso sí, que pasar de media hora
produce el efecto contrario: la persona despierta aturdida.

\pa No todos están convencidos. Un gerente de ventas, que pidió no ser nombrado, opina que «la siesta es un premio a la
pereza» y que sus vendedores «deberían dormir en la casa».

\pa Las empresas que la han adoptado reportan menos errores en las tareas de la tarde, aunque reconocen que sus datos son
internos y no han sido revisados por investigadores externos.

\pa Quizás el debate no sea sobre la siesta en sí, sino sobre una idea más antigua: que trabajar mucho y trabajar bien son
la misma cosa.
\end{texto}
\begin{ejercicios}
\item La cita del médico Andrés Vélez sirve en el texto para
\opcionesl{apoyar con conocimiento experto los beneficios de la siesta corta.}{refutar la siesta.}{dar un testimonio personal.}{presentar la postura del gerente.}
\sol{A. (3.1)}
\item ¿Qué debilidad reconocen las empresas sobre sus propios datos?
\opcionesl{Que son muy antiguos.}{Que son internos y no han sido revisados por investigadores externos.}{Que muestran más errores.}{Que los hizo el gerente.}
\sol{B. (3.2)}
\item La postura del gerente de ventas es
\opcionesl{a favor de la siesta.}{en contra de la siesta.}{neutral.}{la misma del médico.}
\sol{B. (3.2)}
\item Según el último párrafo, el autor sugiere que el debate de fondo es sobre
\opcionesl{el horario de las empresas.}{la creencia de que trabajar mucho equivale a trabajar bien.}{el precio del café.}{la salud de los médicos.}
\sol{B. (3.1)}
\end{ejercicios}
